\begin{exercise}{111}
Show that a cocartesian fibration $p\colon X\to Y$ whose fibres are Joyal-equivalent
to $\Delta^0$ is a trivial fibration.

\begin{proof}
Consider a boundary lifting problem over
$\sigma\colon\Delta^n\to Y$, and pull $p$ back to
$p_\sigma\colon X_\sigma=X\times_Y\Delta^n\to\Delta^n$.
This is again a cocartesian fibration.
Indeed, inner-horn fillers pull back, and a chosen cocartesian lift
in $X$ pairs with the prescribed edge in $\Delta^n$.
Its relative outer-horn fillers also pair with the prescribed base
simplex, so Definition~3.1.1 verifies that it remains cocartesian.
The fibres are the original $X_{\sigma(i)}$.

Since $\Delta^n$ is an $\infty$-category and $p_\sigma$ is an inner
fibration, $X_\sigma$ is an $\infty$-category.
Regard $p_\sigma$ as a morphism from this cocartesian fibration to
$\mathrm{id}_{\Delta^n}$ over $\Delta^n$.
It preserves cocartesian edges, since every edge is cocartesian for
the identity fibration.
Every fibre map $X_{\sigma(i)}\to\Delta^0$ is a Joyal equivalence,
so Theorem~3.1.27 makes $p_\sigma$ a Joyal equivalence.
Its fibres are nonempty, hence it is surjective on vertices.
Exercise~\ref{ex:104} therefore makes $p_\sigma$ a trivial fibration.

The given boundary map lifts to $X_\sigma$ and extends to $\Delta^n$;
composing with $X_\sigma\to X$ solves the original problem.
This works for every $n\ge0$, so $p$ is a trivial fibration.
\end{proof}
\end{exercise}
